% Diese datei bewahrt was ein Beischpiel für die Erklerungen über ein Deutschsprachethema

\documentclass[a4paper,11pt]{article}

\usepackage[a4paper,margin=2cm]{geometry}
\usepackage[T1]{fontenc}
\usepackage[utf8]{inputenc}
\usepackage[ngerman,english]{babel}
\usepackage{lmodern}
\usepackage{microtype}
\usepackage[table]{xcolor}
\usepackage{parskip}
\usepackage{enumitem}
\usepackage[most]{tcolorbox}
\usepackage{titlesec}
\usepackage[hidelinks]{hyperref}

\definecolor{HeaderBlueDark}{RGB}{70,110,170}
\definecolor{RowBlueLight}{RGB}{235,243,255}
\definecolor{RowBlueDark}{RGB}{220,233,250}
\definecolor{SoftGray}{RGB}{90,90,90}

\setlength{\parindent}{0pt}
\setlist[itemize]{leftmargin=1.4em,itemsep=0.35em,topsep=0.35em}
\linespread{1.06}

\titleformat{\section}[block]
{\Large\bfseries\color{HeaderBlueDark}}
{}{0pt}{}
\titlespacing{\section}{0pt}{1.3em}{0.8em}

\titleformat{\subsection}[block]
{\large\bfseries\color{HeaderBlueDark}}
{}{0pt}{}
\titlespacing{\subsection}{0pt}{1.0em}{0.6em}

\newtcolorbox{exbox}{enhanced,breakable,colback=RowBlueLight,colframe=RowBlueDark,boxrule=0.4pt,arc=1.5mm,left=1.2mm,right=1.2mm,top=1mm,bottom=1mm,title={Beispiele \textnormal{(Examples)}},fonttitle=\bfseries,colbacktitle=RowBlueDark,coltitle=black}

\NewDocumentEnvironment{grammarentry}{mm+b}{%
    \begin{tcolorbox}[enhanced,breakable,colback=white,colframe=HeaderBlueDark,boxrule=0.6pt,arc=1.5mm,left=1.5mm,right=1.5mm,top=1mm,bottom=1mm,colbacktitle=HeaderBlueDark,coltitle=white,fonttitle=\bfseries,title={#1\hfill\normalfont\footnotesize #2}]
    #3
    \end{tcolorbox}
}{}

\newcommand{\GermanText}[1]{\textbf{Deutsch: }#1\par}
\newcommand{\EnglishText}[1]{{\small\color{SoftGray}\textbf{English: }#1\par}}
\newcommand{\ExampleItem}[2]{\item \textbf{#1}\\{\small\color{SoftGray}#2}}

\usepackage{array}
\usepackage{longtable}
\title{Die Präposition unter (The preposition under / below / among)}
\date{}
\begin{document}
\maketitle
\tableofcontents
\newpage
\section{Grundbedeutung (Basic meaning)}
\begin{grammarentry}{unter}{under / below / among}
\GermanText{\textbf{unter} bedeutet räumlich: unterhalb von etwas oder unter mehreren Personen.}
\EnglishText{\textbf{unter} in spatial use means: under something or among several people.}
\end{grammarentry}
\section{Wo? + Dativ (Location: Where? + dative)}
\begin{grammarentry}{unter + Dativ}{location}
\GermanText{Der Ausdruck bezeichnet die Position, an der sich jemand oder etwas befindet oder eine Handlung abspielt.}
\EnglishText{The expression describes the location where somebody or something is or where an action takes place.}
\end{grammarentry}
\begin{exbox}
\begin{itemize}
\ExampleItem{Die Katze schläft unter dem Tisch.}{The cat is sleeping under the table.}
\end{itemize}
\end{exbox}
\section{Wohin? + Akkusativ (Destination: Where to? + accusative)}
\begin{grammarentry}{unter + Akkusativ}{destination}
\GermanText{Der Ausdruck bezeichnet den Zielort einer Bewegung oder Platzierung.}
\EnglishText{The expression describes the destination of a movement or placement.}
\end{grammarentry}
\begin{exbox}
\begin{itemize}
\ExampleItem{Die Katze kriecht unter den Tisch.}{The cat crawls under the table.}
\end{itemize}
\end{exbox}
\section{Übertragene und feste Verwendungen (Figurative and fixed uses)}
\begin{grammarentry}{Weitere Bedeutungen}{other uses}
\GermanText{\textbf{unter Druck stehen} (Zustand) und \textbf{jemanden unter Druck setzen} (Veränderung) stehen einander gegenüber. \textbf{unter der Woche} ist ein zeitlicher Ausdruck mit Dativ.}
\EnglishText{Compare \textbf{unter Druck stehen} (be under pressure) with \textbf{jemanden unter Druck setzen} (put someone under pressure). \textbf{unter der Woche} is a temporal dative phrase.}
\end{grammarentry}
\begin{exbox}
\begin{itemize}
\ExampleItem{Unter den Gästen waren viele Studierende.}{Many students were among the guests.}
\ExampleItem{Sie möchte wieder unter die Leute gehen.}{She wants to go out and meet people again.}
\ExampleItem{Unter der Woche habe ich wenig Zeit.}{I have little time during the week.}
\end{itemize}
\end{exbox}
\section{Vergleich und häufiger Fehler (Comparison and common error)}
\begin{grammarentry}{Ort oder Ziel?}{location or destination}
\GermanText{\textbf{Die Katze schläft unter dem Tisch.} (Dativ, Wo?) im Gegensatz zu \textbf{Die Katze kriecht unter den Tisch.} (Akkusativ, Wohin?). Nicht jedes Bewegungsverb verlangt den Akkusativ.}
\EnglishText{\textbf{Die Katze schläft unter dem Tisch.} (dative, location) contrasts with \textbf{Die Katze kriecht unter den Tisch.} (accusative, destination). Movement verbs do not automatically require accusative.}
\end{grammarentry}
\section{Zusammenfassung (Summary)}
\begin{grammarentry}{Merksatz}{key rule}
\GermanText{Räumlich: \textbf{Wo? = Dativ, Wohin? = Akkusativ}. Die feste Rektion muss getrennt gelernt werden.}
\EnglishText{Spatial use: \textbf{Where? = dative, Where to? = accusative}. Learn fixed case government separately.}
\end{grammarentry}
\end{document}
